\section{Related Work}
\label{sec:related}

\subsection{Community structure in social networks}
Community detection partitions a graph into groups that are densely connected inside
and sparsely connected to each other. The usual quality measure is
modularity~\citep{newman2006modularity}, which compares the share of edges inside
communities with the share expected in a random graph that has the same degrees. The
Louvain method~\citep{blondel2008louvain} optimises modularity greedily. Leiden
\citep{traag2019leiden} refines it and guarantees that every community is connected;
we use Leiden. Detection algorithms always return a partition, even on random graphs,
so it is standard to compare against a degree-preserving randomised
graph~\citep{maslov2002specificity}. Separate runs of the same stochastic algorithm
also disagree, and \citet{riolo2020consistency} use that run-to-run variation to
measure how consistently each node is assigned. We use it as a structure-only
baseline.

\subsection{Node metadata versus network structure}
Users carry metadata (biographies, posts, declared interests) alongside their edges.
\citet{newman2016annotated} show that metadata and structural communities can agree
or disagree, and that how well the metadata predicts the community is itself
informative. This is the closest precedent for our content view, which asks what
community a node's text predicts. Community outlier detection~\citep{gao2010outliers}
flags nodes whose attributes are far from the rest of their community. We use a simple
version of that idea (distance from the community's text centroid) as a baseline.

\subsection{Bridges and ambiguous membership}
Several lines of work suggest that nodes between communities matter more for spreading
than their degree implies. The participation coefficient of
\citet{guimera2005functional} measures how evenly a node's links are spread across
communities. \citet{nepusz2008fuzzy} define \emph{bridgeness} from fuzzy (soft)
community memberships, so a node with ambiguous membership is a bridge. In epidemic
models on community-structured networks, \citet{salathe2010dynamics} found that
immunising community bridges can beat immunising high-degree nodes when community
structure is strong. On Telegram, \citet{shah2024telegram} identified channels that
connect misinformation communities and proposed them as intervention targets. All of
these define bridges or ambiguity from structure alone.

\subsection{Spread models and influence blocking}
The Independent Cascade (IC) model~\citep{kempe2003maximizing} spreads a message along
each edge independently with some probability. Its \emph{weighted cascade} variant
sets that probability to $1/\deg(v)$ for the receiving node $v$. Influence Blocking
Maximization (IBM) chooses nodes or links to block so that a harmful cascade
reaches as few users as possible~\citep{budak2011limiting}. The problem is NP-hard, and
unlike influence maximisation its objective is not supermodular, so greedy algorithms
come with no general approximation guarantee~\citep{xie2025blocking}. Community-aware
variants choose blockers using the community partition~\citep{ghoshal2021community}.
Most IBM work assumes the source of the misinformation is known; \citet{chen2021uncertain}
treat the source as uncertain, which is the setting we adopt. Strong structural
baselines include PageRank, betweenness, adaptive degree and Collective
Influence~\citep{morone2015influence}. Which of these wins depends on how easily the
message is transmitted~\citep{hebert2013global}, so we sweep the transmission
probability rather than fix one value.

Spread on real platforms also depends on content. \citet{delvicario2016spreading}
document selective exposure: users mostly share content that agrees with their
group's views. This motivates our content-aware spread model, in which the chance of
passing a message on depends on how closely the receiver's content matches the
community the message came from.

\subsection{Self-supervised graph embeddings}
The original plan was to embed every node from text and topology together using
self-supervised graph neural networks, then cluster the embedding. We trained three
objective families on a shared encoder: BGRL (bootstrapping, no
negatives)~\citep{thakoor2022bgrl}, Deep Graph Infomax (mutual
information)~\citep{velickovic2019dgi} and MVGRL (adjacency view against diffusion
view)~\citep{hassani2020mvgrl}. \Cref{sec:method-evolution} explains why this design
was replaced.

\subsection{Active learning with clustering ambiguity}
In graph active learning, AGE~\citep{cai2017age} picks nodes to label by combining
prediction entropy, centrality and distance to cluster centres. So ``uncertainty plus
centrality'' is already an established design, though its uncertainty is about a
node's \emph{class label}. AAS~\citep{sani2026aas}, built for animal
re-identification, takes a different route. It runs two complementary clusterings
(density-based and nearest-neighbour) on the same embedding, builds an overlap graph
between their clusters, and treats connected regions of partial overlap as
\emph{ambiguous regions}. It then samples must-link and cannot-link pair queries from
those regions and refines the clustering with the oracle's answers. AAS is the
starting point of this project.

\subsection{From related work to the current idea}
\label{sec:related-gap}
Three observations shaped the current design.
\begin{enumerate}
  \item \textbf{Both halves of the original idea already exist.} Community-aware
        blocking~\citep{ghoshal2021community} and uncertainty combined with
        centrality~\citep{cai2017age} are both published. A defensible contribution has
        to be narrower.
  \item \textbf{Existing ambiguity is structural.} Bridgeness, participation and
        consensus ambiguity~\citep{nepusz2008fuzzy,guimera2005functional,riolo2020consistency}
        are all computed from the graph. None of them can see a node whose edges place
        it in one community and whose content places it in another.
  \item \textbf{AAS-style disagreement needs separate views.} Two clusterings of
        one embedding disagree partly because of the algorithms themselves. If one
        view is derived from the other (for example, a content clustering built from a
        graph embedding), their disagreement cannot be read as a property of the node.
\end{enumerate}
We therefore keep the AAS principle (look where complementary views disagree) and
keep the inputs of the two views separate: topology from edges only, content from text
only. The separation is in the inputs, not complete: the content classifier is trained
to predict the topology communities, so its targets come from the graph. It never
receives edges or graph-derived features, and cross-fitting ensures that no node's own
label informs its prediction. The working claim is that \emph{metadata-aware membership ambiguity}
identifies containment targets that purely structural bridge and centrality scores
miss, and the experiments are designed to find out when, if ever, that claim holds.
